% Options for packages loaded elsewhere
% Options for packages loaded elsewhere
\PassOptionsToPackage{unicode}{hyperref}
\PassOptionsToPackage{hyphens}{url}
\PassOptionsToPackage{dvipsnames,svgnames,x11names}{xcolor}
%
\documentclass[
  italian,
  10pt,
  a4paper,
  DIV=11,
  numbers=noendperiod]{scrartcl}
\usepackage{xcolor}
\usepackage[margin=2.2cm]{geometry}
\usepackage{amsmath,amssymb}
\setcounter{secnumdepth}{-\maxdimen} % remove section numbering
\usepackage{iftex}
\ifPDFTeX
  \usepackage[T1]{fontenc}
  \usepackage[utf8]{inputenc}
  \usepackage{textcomp} % provide euro and other symbols
\else % if luatex or xetex
  \usepackage{unicode-math} % this also loads fontspec
  \defaultfontfeatures{Scale=MatchLowercase}
  \defaultfontfeatures[\rmfamily]{Ligatures=TeX,Scale=1}
\fi
\usepackage{lmodern}
\ifPDFTeX\else
  % xetex/luatex font selection
\fi
% Use upquote if available, for straight quotes in verbatim environments
\IfFileExists{upquote.sty}{\usepackage{upquote}}{}
\IfFileExists{microtype.sty}{% use microtype if available
  \usepackage[]{microtype}
  \UseMicrotypeSet[protrusion]{basicmath} % disable protrusion for tt fonts
}{}
\makeatletter
\@ifundefined{KOMAClassName}{% if non-KOMA class
  \IfFileExists{parskip.sty}{%
    \usepackage{parskip}
  }{% else
    \setlength{\parindent}{0pt}
    \setlength{\parskip}{6pt plus 2pt minus 1pt}}
}{% if KOMA class
  \KOMAoptions{parskip=half}}
\makeatother
% Make \paragraph and \subparagraph free-standing
\makeatletter
\ifx\paragraph\undefined\else
  \let\oldparagraph\paragraph
  \renewcommand{\paragraph}{
    \@ifstar
      \xxxParagraphStar
      \xxxParagraphNoStar
  }
  \newcommand{\xxxParagraphStar}[1]{\oldparagraph*{#1}\mbox{}}
  \newcommand{\xxxParagraphNoStar}[1]{\oldparagraph{#1}\mbox{}}
\fi
\ifx\subparagraph\undefined\else
  \let\oldsubparagraph\subparagraph
  \renewcommand{\subparagraph}{
    \@ifstar
      \xxxSubParagraphStar
      \xxxSubParagraphNoStar
  }
  \newcommand{\xxxSubParagraphStar}[1]{\oldsubparagraph*{#1}\mbox{}}
  \newcommand{\xxxSubParagraphNoStar}[1]{\oldsubparagraph{#1}\mbox{}}
\fi
\makeatother


\usepackage{longtable,booktabs,array}
\newcounter{none} % for unnumbered tables
\usepackage{calc} % for calculating minipage widths
% Correct order of tables after \paragraph or \subparagraph
\usepackage{etoolbox}
\makeatletter
\patchcmd\longtable{\par}{\if@noskipsec\mbox{}\fi\par}{}{}
\makeatother
% Allow footnotes in longtable head/foot
\IfFileExists{footnotehyper.sty}{\usepackage{footnotehyper}}{\usepackage{footnote}}
\makesavenoteenv{longtable}
\usepackage{graphicx}
\makeatletter
\newsavebox\pandoc@box
\newcommand*\pandocbounded[1]{% scales image to fit in text height/width
  \sbox\pandoc@box{#1}%
  \Gscale@div\@tempa{\textheight}{\dimexpr\ht\pandoc@box+\dp\pandoc@box\relax}%
  \Gscale@div\@tempb{\linewidth}{\wd\pandoc@box}%
  \ifdim\@tempb\p@<\@tempa\p@\let\@tempa\@tempb\fi% select the smaller of both
  \ifdim\@tempa\p@<\p@\scalebox{\@tempa}{\usebox\pandoc@box}%
  \else\usebox{\pandoc@box}%
  \fi%
}
% Set default figure placement to htbp
\def\fps@figure{htbp}
\makeatother



\ifLuaTeX
\usepackage[bidi=basic,shorthands=off]{babel}
\else
\usepackage[bidi=default,shorthands=off]{babel}
\fi
\ifLuaTeX
  \usepackage{selnolig} % disable illegal ligatures
\fi


\setlength{\emergencystretch}{3em} % prevent overfull lines

\providecommand{\tightlist}{%
  \setlength{\itemsep}{0pt}\setlength{\parskip}{0pt}}



 


\KOMAoption{captions}{tableheading}
\makeatletter
\@ifpackageloaded{tcolorbox}{}{\usepackage[skins,breakable]{tcolorbox}}
\@ifpackageloaded{fontawesome5}{}{\usepackage{fontawesome5}}
\definecolor{quarto-callout-color}{HTML}{909090}
\definecolor{quarto-callout-note-color}{HTML}{0758E5}
\definecolor{quarto-callout-important-color}{HTML}{CC1914}
\definecolor{quarto-callout-warning-color}{HTML}{EB9113}
\definecolor{quarto-callout-tip-color}{HTML}{00A047}
\definecolor{quarto-callout-caution-color}{HTML}{FC5300}
\definecolor{quarto-callout-color-frame}{HTML}{acacac}
\definecolor{quarto-callout-note-color-frame}{HTML}{4582ec}
\definecolor{quarto-callout-important-color-frame}{HTML}{d9534f}
\definecolor{quarto-callout-warning-color-frame}{HTML}{f0ad4e}
\definecolor{quarto-callout-tip-color-frame}{HTML}{02b875}
\definecolor{quarto-callout-caution-color-frame}{HTML}{fd7e14}
\makeatother
\makeatletter
\@ifpackageloaded{caption}{}{\usepackage{caption}}
\AtBeginDocument{%
\ifdefined\contentsname
  \renewcommand*\contentsname{Indice}
\else
  \newcommand\contentsname{Indice}
\fi
\ifdefined\listfigurename
  \renewcommand*\listfigurename{Elenco delle Figure}
\else
  \newcommand\listfigurename{Elenco delle Figure}
\fi
\ifdefined\listtablename
  \renewcommand*\listtablename{Elenco delle Tabelle}
\else
  \newcommand\listtablename{Elenco delle Tabelle}
\fi
\ifdefined\figurename
  \renewcommand*\figurename{Figura}
\else
  \newcommand\figurename{Figura}
\fi
\ifdefined\tablename
  \renewcommand*\tablename{Tabella}
\else
  \newcommand\tablename{Tabella}
\fi
}
\@ifpackageloaded{float}{}{\usepackage{float}}
\floatstyle{ruled}
\@ifundefined{c@chapter}{\newfloat{codelisting}{h}{lop}}{\newfloat{codelisting}{h}{lop}[chapter]}
\floatname{codelisting}{Lista}
\newcommand*\listoflistings{\listof{codelisting}{Elenco degli Elenchi}}
\makeatother
\makeatletter
\makeatother
\makeatletter
\@ifpackageloaded{caption}{}{\usepackage{caption}}
\@ifpackageloaded{subcaption}{}{\usepackage{subcaption}}
\makeatother
\usepackage{bookmark}
\IfFileExists{xurl.sty}{\usepackage{xurl}}{} % add URL line breaks if available
\urlstyle{same}
% fallback for those not using the hyperref driver hyperxmp:
\makeatletter
\@ifundefined{xmpquote}{\newcommand{\xmpquote}[1]{#1}}{}
\makeatother
\hypersetup{
  pdftitle={Curriculum Vitae},
  pdflang={it},
  colorlinks=true,
  linkcolor={blue},
  filecolor={Maroon},
  citecolor={Blue},
  urlcolor={Blue},
  pdfcreator={LaTeX via pandoc}}


\title{Curriculum Vitae}
\usepackage{etoolbox}
\makeatletter
\providecommand{\subtitle}[1]{% add subtitle to \maketitle
  \apptocmd{\@title}{\par {\large #1 \par}}{}{}
}
\makeatother
\subtitle{Giovanni Castellazzi --- Professore associato, Università di
Bologna}
\author{}
\date{4 novembre 2026}
\begin{document}
\maketitle

\renewcommand*\contentsname{Indice}
{
\setcounter{tocdepth}{2}
\tableofcontents
}

\begin{tcolorbox}[enhanced jigsaw, arc=.35mm, bottomrule=.15mm, breakable, colback=white, colframe=quarto-callout-note-color-frame, left=2mm, leftrule=.75mm, opacityback=0, rightrule=.15mm, toprule=.15mm]
\begin{minipage}[t]{5.5mm}
\textcolor{quarto-callout-note-color}{\faInfo}
\end{minipage}%
\begin{minipage}[t]{\textwidth - 5.5mm}

Profili: \href{https://orcid.org/0000-0001-5870-5474}{ORCID} ·
\href{https://scholar.google.com/citations?user=pPopX2QAAAAJ&hl=it}{Google
Scholar} ·
\href{https://gcastellazzi.github.io/publications.html}{Pubblicazioni}

\end{minipage}%
\end{tcolorbox}

\subsection{Posizione attuale}\label{posizione-attuale}

\textbf{Professore associato} --- Dipartimento di Ingegneria Civile,
Chimica, Ambientale e dei Materiali (DICAM), Università di Bologna.
\emph{Dal 2021.}

\textbf{Presidente} --- AICO, Associazione Italiana Compositi per le
applicazioni nell'ingegneria civile. \emph{Dal 2023.}
\href{https://www.aico-compositi.it}{aico-compositi.it}

\textbf{Membro del collegio docenti} --- Dottorato in \emph{Ingegneria e
tecnologia dell'informazione per il monitoraggio strutturale e
ambientale e la gestione dei rischi} (EIT4SEMM), Università di Bologna.
\emph{Dal 2022.}

\textbf{Ingegnere iscritto all'Albo} --- Ordine degli Ingegneri della
Provincia di Modena. \emph{Dal 2002.}

\subsection{Abilitazione scientifica
nazionale}\label{abilitazione-scientifica-nazionale}

{\def\LTcaptype{none} % do not increment counter
\begin{longtable}[]{@{}
  >{\raggedright\arraybackslash}p{(\linewidth - 2\tabcolsep) * \real{0.2200}}
  >{\raggedright\arraybackslash}p{(\linewidth - 2\tabcolsep) * \real{0.7800}}@{}}
\toprule\noalign{}
\begin{minipage}[b]{\linewidth}\raggedright
Periodo
\end{minipage} & \begin{minipage}[b]{\linewidth}\raggedright
Abilitazione
\end{minipage} \\
\midrule\noalign{}
\endhead
\bottomrule\noalign{}
\endlastfoot
2025--2035 & \textbf{Professore ordinario} --- CEAR-06/A, Meccanica dei
solidi e delle strutture \\
2017--2028 & Professore associato --- SSD 08/B2, Scienza delle
costruzioni \\
\end{longtable}
}

\subsection{Formazione}\label{formazione}

{\def\LTcaptype{none} % do not increment counter
\begin{longtable}[]{@{}
  >{\raggedright\arraybackslash}p{(\linewidth - 2\tabcolsep) * \real{0.1600}}
  >{\raggedright\arraybackslash}p{(\linewidth - 2\tabcolsep) * \real{0.8400}}@{}}
\toprule\noalign{}
\begin{minipage}[b]{\linewidth}\raggedright
Periodo
\end{minipage} & \begin{minipage}[b]{\linewidth}\raggedright
\end{minipage} \\
\midrule\noalign{}
\endhead
\bottomrule\noalign{}
\endlastfoot
2009--2010 & \textbf{Assegnista post-dottorato}, University of
California San Diego, Jacobs School of Engineering --- \emph{Multi-level
schemes for explicit dynamics to track elastic/acoustic waves} \\
2008--2009 & \textbf{Assegnista post-dottorato}, University of
California San Diego, Jacobs School of Engineering ---
\emph{Assumed-strain elements for incompressible media and plates} \\
2004--2007 & \textbf{Dottorato in Meccanica delle strutture}, Università
di Bologna \\
2001 & \textbf{Laurea in Ingegneria Civile} (quinquennale, indirizzo
Strutture), Università di Bologna \\
\end{longtable}
}

\subsection{Posizioni accademiche}\label{posizioni-accademiche}

{\def\LTcaptype{none} % do not increment counter
\begin{longtable}[]{@{}
  >{\raggedright\arraybackslash}p{(\linewidth - 2\tabcolsep) * \real{0.2200}}
  >{\raggedright\arraybackslash}p{(\linewidth - 2\tabcolsep) * \real{0.7800}}@{}}
\toprule\noalign{}
\begin{minipage}[b]{\linewidth}\raggedright
Periodo
\end{minipage} & \begin{minipage}[b]{\linewidth}\raggedright
Posizione
\end{minipage} \\
\midrule\noalign{}
\endhead
\bottomrule\noalign{}
\endlastfoot
2021--oggi & Professore associato, DICAM, Università di Bologna \\
2018--2021 & Ricercatore a tempo determinato di tipo B (RTDb), DICAM,
Università di Bologna \\
2001--2018 & Tecnico laureato di supporto al laboratorio, DICAM (già
DISTART), Università di Bologna \\
\end{longtable}
}

\subsection{Incarichi istituzionali}\label{incarichi-istituzionali}

\subsubsection{Corsi di studio e
commissioni}\label{corsi-di-studio-e-commissioni}

{\def\LTcaptype{none} % do not increment counter
\begin{longtable}[]{@{}
  >{\raggedright\arraybackslash}p{(\linewidth - 2\tabcolsep) * \real{0.2200}}
  >{\raggedright\arraybackslash}p{(\linewidth - 2\tabcolsep) * \real{0.7800}}@{}}
\toprule\noalign{}
\begin{minipage}[b]{\linewidth}\raggedright
Periodo
\end{minipage} & \begin{minipage}[b]{\linewidth}\raggedright
Incarico
\end{minipage} \\
\midrule\noalign{}
\endhead
\bottomrule\noalign{}
\endlastfoot
2021--2023 & Membro della commissione di accesso, LM in
\href{https://corsi.unibo.it/2cycle/HistoricBuildings}{Historic Building
Rehabilitation} \\
2021--2022 & Membro della commissione di ammissione al dottorato in
Ingegneria Civile, Chimica, Ambientale e dei Materiali (37° ciclo,
membro effettivo) \\
2020--2021 & Membro della commissione di ammissione al dottorato in
Ingegneria Civile, Chimica, Ambientale e dei Materiali (36° ciclo,
membro supplente) \\
2019--2023 & Segretario del corso di laurea magistrale in
\href{https://corsi.unibo.it/2cycle/CivilEngineering}{Ingegneria
Civile} \\
2019--2023 & Membro della
\href{https://corsi.unibo.it/2cycle/CivilEngineering/how-to-enrol}{commissione
di accesso}, LM in Ingegneria Civile \\
2019--2020 & Membro della commissione di Ateneo per gli scarichi
inventariali, Università di Bologna \\
\end{longtable}
}

\subsubsection{Società scientifiche}\label{societuxe0-scientifiche}

{\def\LTcaptype{none} % do not increment counter
\begin{longtable}[]{@{}
  >{\raggedright\arraybackslash}p{(\linewidth - 2\tabcolsep) * \real{0.2200}}
  >{\raggedright\arraybackslash}p{(\linewidth - 2\tabcolsep) * \real{0.7800}}@{}}
\toprule\noalign{}
\begin{minipage}[b]{\linewidth}\raggedright
Periodo
\end{minipage} & \begin{minipage}[b]{\linewidth}\raggedright
Società
\end{minipage} \\
\midrule\noalign{}
\endhead
\bottomrule\noalign{}
\endlastfoot
2023--oggi & \textbf{Presidente}, AICO --- Associazione Italiana
Compositi per le applicazioni nell'ingegneria civile \\
2019--2023 & Vicepresidente, AICO \\
2014--oggi & Socio, AICO ---
\href{https://www.aico-compositi.it}{aico-compositi.it} \\
2019--oggi & Staff member, RILEM --- International Union of Laboratories
and Experts in Construction Materials, Systems and Structures ---
\href{https://www.rilem.net}{rilem.net} \\
2018--oggi & Socio, EUROMECH --- European Mechanics Society \\
2018--oggi & Socio, SISCO --- Associazione Italiana di Scienza delle
Costruzioni \\
2011--oggi & Socio, AIMETA --- Associazione Italiana di Meccanica
Teorica e Applicata --- \href{https://www.aimeta.it}{aimeta.it} \\
2009--2011 & Socio, ISSNAF --- Italian Scientists and Scholars in North
America Foundation \\
2012 & Socio, SPIE --- International Society for Optics and Photonics \\
\end{longtable}
}

\subsection{Supervisione}\label{supervisione}

\subsubsection{Tesi di dottorato}\label{tesi-di-dottorato}

\textbf{2022--2025 · Relatore} --- \emph{Innovative eco-sustainable
materials for the energy and seismic improvement of civil buildings.}
Dottorato EIT4SEMM. Candidato: \textbf{Nicolò Lo Presti}.

\textbf{2015--2018 · Correlatore} --- \emph{Micro-modelling of masonry
structures.} Dottorato in Ingegneria Civile, Chimica, Ambientale e dei
Materiali. Candidato: \textbf{Antonio Maria D'Altri}.

\subsubsection{Assegni di ricerca}\label{assegni-di-ricerca}

\textbf{2020--2021 · Responsabile} --- Assegno di ricerca su
\emph{Existing Building Information Modeling per la gestione
dell'intervento sul costruito esistente}, finanziato nell'ambito del
progetto eBIM (POR-FESR, Regione Emilia-Romagna). Assegnista:
\textbf{Nicolò Lo Presti}.

\textbf{Set--dic 2021 · Responsabile} --- Visiting researcher dalla
Ivane Javakhishvili Tbilisi State University, Georgia, su
\emph{Effective modeling of complex unreinforced masonry buildings}.
Ricercatrice: \textbf{Dr.~Nino Tabatadze}.

\subsection{Didattica}\label{didattica}

Università di Bologna, anno accademico 2026/2027.

{\def\LTcaptype{none} % do not increment counter
\begin{longtable}[]{@{}
  >{\raggedright\arraybackslash}p{(\linewidth - 4\tabcolsep) * \real{0.1000}}
  >{\raggedright\arraybackslash}p{(\linewidth - 4\tabcolsep) * \real{0.4500}}
  >{\raggedright\arraybackslash}p{(\linewidth - 4\tabcolsep) * \real{0.4500}}@{}}
\toprule\noalign{}
\begin{minipage}[b]{\linewidth}\raggedright
Codice
\end{minipage} & \begin{minipage}[b]{\linewidth}\raggedright
Insegnamento
\end{minipage} & \begin{minipage}[b]{\linewidth}\raggedright
Corso di studio
\end{minipage} \\
\midrule\noalign{}
\endhead
\bottomrule\noalign{}
\endlastfoot
72782 & \textbf{Mechanics of Historical Masonry Structures} --- 6 CFU &
LM in Ingegneria Edile--Architettura · LM in Ingegneria Civile \\
72772 & \textbf{Computational Structural Mechanics M} \emph{(modulo 2)}
& LM in Ingegneria Civile \\
B4954 & \textbf{Scienza delle Costruzioni T con Laboratorio}
\emph{(modulo 2)} & L in Ingegneria--Architettura \\
\end{longtable}
}

Insegnamenti e appelli d'esame sono elencati sulla mia
\href{https://www.unibo.it/sitoweb/giovanni.castellazzi/teachings}{pagina
didattica dell'Università di Bologna}.

\begin{tcolorbox}[enhanced jigsaw, arc=.35mm, bottomrule=.15mm, breakable, colback=white, colframe=quarto-callout-note-color-frame, left=2mm, leftrule=.75mm, opacityback=0, rightrule=.15mm, toprule=.15mm]
\begin{minipage}[t]{5.5mm}
\textcolor{quarto-callout-note-color}{\faInfo}
\end{minipage}%
\begin{minipage}[t]{\textwidth - 5.5mm}

\emph{Mechanics of Historical Masonry Structures} si avvale di
\href{https://gcastellazzi.github.io/aLoTiA/app/}{aLOTofImaginArches},
uno strumento interattivo scritto per il corso: lo studente carica la
fotografia di un arco reale, ne traccia la geometria e muove la linea
delle pressioni nei suoi tre gradi di libertà, verificando giunto per
giunto se resta dentro la muratura.

\end{minipage}%
\end{tcolorbox}

\subsection{Contatti}\label{contatti}

Dipartimento di Ingegneria Civile, Chimica, Ambientale e dei Materiali
(DICAM) Università di Bologna · Viale del Risorgimento 2, 40136 Bologna
\href{mailto:giovanni.castellazzi@unibo.it}{\nolinkurl{giovanni.castellazzi@unibo.it}}




\end{document}
